\begin{table}[htbp]
\centering
\caption{Kinerja penelusuran pada basis pengetahuan penuh (6 kasus uji, \texttt{top\_\allowbreak k}=5): RAG standar lawan RAG \emph{terkondisi-prediksi}. ``Butuh target'' menunjukkan proporsi kasus yang menuntut dokumen kondisi masa depan. \textbf{Diukur pada konfigurasi penelusuran padat}, yakni sebelum jalur produksi berpindah ke pencocokan leksikal; angka pada tabel ini karena itu tidak sebanding dengan Tabel~\ref{tab:crossfold-retrieval} dan~\ref{tab:realcases}, sehingga dipertahankan hanya sebagai demonstrasi mekanisme.}
\label{tab:kinerja-retrieval}
\small
\begin{tabular}{@{}lrrrrr@{}}
\toprule
\textbf{Mode} & \textbf{Butuh target (\%)} & \textbf{Hit@1 (\%)} & \textbf{Hit@5 (\%)} & \textbf{MRR} & \textbf{nDCG@5} \\
\midrule
RAG standar & 0,0 & 0,0 & 100,0 & 0,225 & 0,160 \\
RAG terkondisi-prediksi & 100,0 & 83,3 & 100,0 & \textbf{0,889} & \textbf{0,384} \\
\bottomrule
\end{tabular}
\end{table}
